%--------------------------------------------------------------------------------
%
%
%--------------------------------------------------------------------------------
\pagebreak
\section*{LDI Load immediate}
\addcontentsline{toc}{section}{LDI Load immediate}

\instructionentry{Syntax}{
  \texttt{LDI[.L/S/U] RegR, val }
}

\instructionentry{Format}{

	The LDI instruction uses the immediate-20 instruction format. \\

   	\begin{flushleft}
	\begin{tikzpicture}[scale=0.7, transform shape]
        
        % \draw[help lines, gray!70, dashed] (0,0) grid(16,1.5);
        
        \node[  tsRectangle, 
                minimum width=16cm,
                minimum height=1cm,
                text width=3cm,
                text centered,
                fill=white!50] (blocks) at (8,1) { };
                
     	\draw[-] (1,0.5) -- (1,1.5);
     	\draw[-] (3,0.5) -- (3,1.5);
     	\draw[-] (5,0.5) -- (5,1.5);
     	\draw[-] (6,0.5) -- (6,1.5);
     	     	
     	\node at (0.5,0.25) {2};
    	\node at (2,0.25) {4};
    	\node at (4,0.25) {4};
    	\node at (5.5,0.25) {2};
   	 	\node at (11,0.25) {20};
   	 	
   	 	\node at (0.5,1) {\small 0};
		\node at (2,1) {\small \OpCodeADDIL};
		\node at (4,1) {\small Reg R};
		\node at (5.5,1) {\small opt};
		\node at (11,1) {\small val};
   	 	
	\end{tikzpicture}
	\end{flushleft}
}

\instructionentry{Description}{

The \texttt{LDI} instruction ... }

\instructionentry{Operation}{

  	\texttt{RegR <- (( val \& 0xFFFFF ) << 12 ); }  (.L option ) \\
  	\texttt{RegR <- (( val \& 0xFFFFF ) << 32 ); }  (.S option ) \\
  	\texttt{RegR <- (( val \& 0xFFF   ) << 52 ); }  (.U option ) \\
} 

\instructionentry{Exceptions}{

	None.
}

\instructionentry{Notes}{

	None. 
}




